\section{Randomized revealing optimization: the exact four-piece curve}
\label{sec:random-common-upper}

The previous section identified the optimal leading online cost for every
fixed input in the revealing-optimization model.  We now compare that cost with the
clairvoyant optimum and maximize over all inputs.  The cap $u$ is fixed, but
the algorithm is not told the processing-time multiset.  This section proves
\cref{thm:random-common-upper}, the rigorous randomized half of the informal
curve theorem \cref{thm:intro-common-upper-curves}.  Its three transition
points are
\begin{equation}
 \uRand{1}=1,
 \qquad
 \uRand{2}=5.048917339522306\ldots,
 \qquad
 \uRand{3}=\frac{25}{4}.
 \label{eq:rand-transition-points}
\end{equation}
Here $\uRand{2}>1$ is the unique root of
$u^3-6u^2+5u-1=0$ above one.

\begin{theorem}[Randomized revealing-optimization curve]
\label{thm:random-common-upper}
For every fixed $u>0$, the optimal randomized size-asymptotic competitive
ratio against an oblivious adversary is
\begin{equation}
 \RrandRO(u)=
 \begin{cases}
  1,&0<u\le\uRand{1},\\[2pt]
  \displaystyle\frac{u^3}{u^2+(u-1)^3},
    &\uRand{1}\le u\le\uRand{2},\\[8pt]
  \displaystyle 1+\frac{1}{2(\sqrt u-1)},
    &\uRand{2}\le u\le\uRand{3},\\[8pt]
  \displaystyle\frac43,&u\ge\uRand{3}.
 \end{cases}
 \label{eq:intro-random-common-upper-curve}
\end{equation}
More precisely, there are randomized nonanticipating algorithms
$\mathcal A_{n,u}$ and numbers $\eps_u(n)\to0$ such that every fixed labeled
vector $p\in[0,u]^n$ satisfies
\begin{equation}
 \EE\ALG_{\mathcal A_{n,u}}(p)
 \le \RrandRO(u)\OPT_u(p)+\eps_u(n)n^2.
 \label{eq:intro-random-common-upper}
\end{equation}
Conversely, for every randomized online algorithm $\mathcal A$ and every
$\eps>0$, there
are arbitrarily large $n$ and a fixed vector $p\in[0,u]^n$ for which
\begin{equation}
 \EE\ALG_{\mathcal A}(p)
 \ge\bigl(\RrandRO(u)-\eps\bigr)\OPT_u(p).
 \label{eq:random-common-upper-lower}
\end{equation}
The curve is continuous and has global maximum
\begin{equation}
 \max_{u>0}\RrandRO(u)
 =\frac{27+6\sqrt3}{23}
 =1.625752384583185\ldots,
 \label{eq:intro-random-common-maximum}
\end{equation}
attained at $u=(3+\sqrt3)/2$.
\end{theorem}

The preceding instance-optimal theorem identifies the best policy on each
input.  Here we compare that value with the clairvoyant minimum kernel and
maximize over all empirical distributions.  A survival-function flattening
argument reduces the outer optimization to binary distributions.  Two
scalar maximizations give the four branches above.  A sublinear sample removes
advance knowledge of the input, and binary random labelings give matching Yao
lower bounds.

\subsection{Offline and stationary fluid coefficients}

Assume initially that $u>1$ and write
\begin{equation}
 s=u-1.
 \label{eq:rcu-s}
\end{equation}
Let $D$ be a finite distribution on $[0,u]$, let $P,P'$ be independent with
law $D$, and use the survival function
\begin{equation}
 S_D(t)=\Prob(P>t),
 \qquad S_D(t)=0\quad(t>u).
 \label{eq:rcu-survival}
\end{equation}
The convention at atoms does not affect any integral below.  Since a
clairvoyant job has effective length $1+\min\{s,P\}$, the layer-cake identity
and \cref{cor:empirical-offline} give, for $P_M\sim D[M]$,
\begin{align}
 \Omega_{\mathsf{RO},u}(D)
 &=\frac12\left(1+\int_0^s S_D(t)^2\,dt\right),
 \label{eq:rcu-offline-fluid}\\
 \OPT_u(M)
 &=n^2\Omega_{\mathsf{RO},u}(D[M])
   +\frac n2\EE\lambda_u(P_M)
 \label{eq:rcu-offline-finite}
\end{align}
for every empirical distribution $D[M]$.

We next specialize the maximum-density module
\cref{lem:maximum-density-module} to a form suited to the cap.  If its
threshold prefix has mass $a$ and first moment $\ell$, put
\begin{equation}
 \tau_D=\frac{1+\ell}{a}.
 \label{eq:rcu-tau-density}
\end{equation}
The common threshold equation and the layer-cake identity give
\begin{equation}
 \EE(\tau_D-P)^+=1,
 \qquad
 \int_0^{\tau_D} S_D(t)\,dt=\tau_D-1.
 \label{eq:rcu-threshold-equation}
\end{equation}

The stationary policy tests all jobs in a uniform random order, processes a
positive prefix outcome immediately after its test, and after the last test
processes the residual jobs in SPT order.  The following formula is the
stationary pair identity specialized to a finite cap.

\begin{lemma}[Stationary survival formula]
\label{lem:rcu-stationary-survival}
The doubled leading coefficient of the stationary policy is
\begin{equation}
 \mathcal T_u(D)=
 \begin{cases}
  \displaystyle \tau_D+\int_{\tau_D}^u S_D(t)^2\,dt,&\tau_D\le u,\\[6pt]
  \tau_D,&\tau_D>u.
 \end{cases}
 \label{eq:rcu-stationary-survival}
\end{equation}
Consequently the better of the raw and stationary policies has leading
ratio
\begin{equation}
 \frac{\min\{u,\mathcal T_u(D)\}}
      {1+\int_0^sS_D(t)^2\,dt}.
 \label{eq:rcu-announced-ratio}
\end{equation}
\end{lemma}

\begin{proof}
Let $h=1-a$ be the residual mass and let $K_T$ be the minimum kernel of two
residual draws, with their subprobability masses left unnormalized.  The
stationary prefix identity \cref{lem:stationary-prefix-cost}, applied first
to empirical laws and then by continuity, gives the leading coefficient
\begin{equation}
 (1+\ell)\left(1-\frac a2\right)+\frac{K_T}{2}.
 \label{eq:rcu-stationary-pair}
\end{equation}
If $\tau_D\le u$, the threshold property gives
\[
 K_T=\tau_D h^2+\int_{\tau_D}^uS_D(t)^2\,dt,
 \qquad
 (1+\ell)\left(1-\frac a2\right)
 =\frac{\tau_D}2(1-h^2).
\]
Their sum is one half of the first line of
\eqref{eq:rcu-stationary-survival}.  If $\tau_D>u$, the whole distribution is
in the prefix and \eqref{eq:rcu-stationary-pair} equals $\tau_D/2$.  Raw
execution has leading coefficient $u/2$, and division by
\eqref{eq:rcu-offline-fluid} proves \eqref{eq:rcu-announced-ratio}.
\end{proof}

This is where the instance-optimal theorem enters the curve calculation.
In \cref{eq:cui-F}, $q=0$ is the raw policy and $q=1$ is precisely the
stationary maximum-density policy.  Hence
\begin{equation}
 2\PhiRO(D)\le\min\{u,\mathcal T_u(D)\},
 \label{eq:rcu-Phi-endpoint-upper}
\end{equation}
and \cref{eq:rcu-announced-ratio} is an upper bound on the
instance-specific ratio $\PhiRO(D)/\Omega_{\mathsf{RO},u}(D)$.  Interior values of $q$ can
strictly improve individual inputs; the point of the next two subsections is
that the two endpoint policies already give the sharp value after maximizing
over all $D$.  The binary lower bounds at the end apply to every policy, not
only to these endpoints.

\subsection{Flattening the survival function}

We now maximize \eqref{eq:rcu-announced-ratio} over arbitrary $D$.  Suppose
first that $1\le\tau_D\le u$ and put
\begin{equation}
 y=\frac{\tau_D-1}{\tau_D}.
 \label{eq:rcu-y}
\end{equation}
By \eqref{eq:rcu-threshold-equation}, $y$ is the average of $S_D$ on
$[0,\tau_D]$.  Since $S_D$ is nonincreasing, its average on every initial
subinterval is at least $y$.  Thus Cauchy--Schwarz gives
\begin{equation}
 \int_0^aS_D(t)^2\,dt\ge ay^2
 \qquad(0\le a\le\tau_D),
 \label{eq:rcu-initial-cs}
\end{equation}
and $S_D(t)\le y$ almost everywhere for $t>\tau_D$.

Replace $S_D$ on $(0,\tau_D)$ by the constant $y$ and leave its tail unchanged.
This remains a nonincreasing survival function, preserves
\eqref{eq:rcu-threshold-equation} and $\mathcal T_u(D)$, and by
\eqref{eq:rcu-initial-cs} can only decrease the offline coefficient.  Hence
it can only increase the ratio.  We may therefore assume
\begin{equation}
 S_D(t)=y\quad(0<t<\tau_D).
 \label{eq:rcu-flat-prefix}
\end{equation}

If $u-1\le\tau_D\le u$, the offline integral lies entirely in the flat part,
and $S_D(t)^2\le y^2$ on the remaining tail.  It follows that
\begin{equation}
 \frac{\min\{u,\mathcal T_u(D)\}}
      {1+\int_0^sS_D(t)^2\,dt}
 \le
 B_u(\tau_D)
 \defeq
 \frac{\tau_D+(u-\tau_D)y^2}{1+(u-1)y^2}.
 \label{eq:rcu-B-first}
\end{equation}
The numerator displayed on the right is at most $u$, so the raw minimum is
inactive.

Now suppose $\tau_D<u-1$.  Put
\begin{equation}
 L=u-1-\tau_D,
 \quad E=\int_{\tau_D}^{u-1}S_D(t)^2\,dt,
 \quad G=\int_{u-1}^uS_D(t)^2\,dt.
 \label{eq:rcu-EG}
\end{equation}
Monotonicity gives
\begin{equation}
 0\le E\le Ly^2,
 \qquad G\le E/L.
 \label{eq:rcu-E-range}
\end{equation}
Consequently
\begin{equation}
 \frac{\min\{u,\mathcal T_u(D)\}}
      {1+\int_0^sS_D(t)^2\,dt}
 \le
 \frac{\tau_D+\dfrac{u-\tau_D}{u-1-\tau_D}E}
      {1+\tau_D y^2+E}.
 \label{eq:rcu-linear-fractional-E}
\end{equation}
For fixed $\tau_D$ this is linear-fractional in $E$, so its derivative has a
constant sign.  Its maximum on $[0,Ly^2]$ occurs at an endpoint.  The two
endpoint values are
\begin{align}
 A(\tau_D)&\defeq\frac{\tau_D}{1+\tau_D y^2}
              =\frac{\tau_D^2}{\tau_D^2-\tau_D+1},
 \label{eq:rcu-A}\\
 B_u(\tau_D)&=\frac{\tau_D+(u-\tau_D)y^2}{1+(u-1)y^2}.
 \label{eq:rcu-B}
\end{align}
They are attained by rectangular survival functions: respectively a binary
distribution on $\{0,\tau_D\}$ and one on $\{0,u\}$.

It remains to cover $\tau_D>u$.  Equation
\eqref{eq:rcu-threshold-equation} and the zero extension of $S_D$ give
$\int_0^uS_D(t)\,dt=\tau_D-1\ge u-1$.  The average of a nonincreasing function on
$[0,u-1]$ is no smaller than its average on $[0,u]$, and hence
\begin{equation}
 \int_0^{u-1}S_D(t)^2\,dt
 \ge (u-1)\left(\frac{\tau_D-1}{u}\right)^2
 \ge\frac{(u-1)^3}{u^2}.
 \label{eq:rcu-tau-large}
\end{equation}
Here the raw policy is chosen, and its ratio is at most $B_u(u)$.
We have proved the reduction
\begin{equation}
 \sup_D\frac{\min\{u,\mathcal T_u(D)\}}
                 {1+\int_0^sS_D(t)^2\,dt}
 \le
 \max\left\{
  \sup_{1\le\tau\le u-1}A(\tau),
  \sup_{1\le\tau\le u}B_u(\tau)
 \right\},
 \label{eq:rcu-binary-reduction}
\end{equation}
where the first supremum is omitted for $u<2$.  In fact equality holds in
\eqref{eq:rcu-binary-reduction}, because the rectangular survival functions
described above attain both families.

\subsection{The two scalar maximizations}

The first family satisfies
\begin{equation}
 A'(\tau)=\frac{\tau(2-\tau)}{(\tau^2-\tau+1)^2}.
 \label{eq:rcu-A-derivative}
\end{equation}
Thus its maximum is attained at $\tau=\min\{2,u-1\}$ and equals $4/3$ as
soon as $u\ge3$.

For the second family, substitution of \eqref{eq:rcu-y} gives
\begin{equation}
 B_u(\tau)=
 \frac{\tau^2u+2\tau^2-2\tau u-\tau+u}
      {\tau^2u-2\tau u+2\tau+u-1}.
 \label{eq:rcu-B-rational}
\end{equation}
The sign of its derivative is the sign of
\begin{equation}
 g_u(\tau)=(2\tau-1)^2-u(\tau-1)^2.
 \label{eq:rcu-B-derivative}
\end{equation}
For $u\le4$ this is positive on $[1,\infty)$.  For $u>4$, its unique root
above one is
\begin{equation}
 \tau_*(u)=\frac{u-2+\sqrt u}{u-4}
           =\frac{\sqrt u-1}{\sqrt u-2},
 \label{eq:rcu-tau-star}
\end{equation}
and $B_u$ increases before this point and decreases afterwards.  The
identity $\tau_*(u)=u$ is equivalent to
\begin{equation}
 u^3-6u^2+5u-1=0.
 \label{eq:rcu-transition-cubic}
\end{equation}
The cubic is strictly increasing beyond its last critical point and changes
sign there exactly once; its unique root above one is $\uRand{2}$ from
\eqref{eq:rand-transition-points}.  Hence
\begin{align}
 \max_{1\le\tau\le u}B_u(\tau)
 &=B_u(u)=\frac{u^3}{u^3-2u^2+3u-1}
       =\frac{u^3}{u^2+(u-1)^3},&&u\le\uRand{2},
 \label{eq:rcu-B-at-u}\\
 &=B_u(\tau_*(u))
       =1+\frac1{2(\sqrt u-1)},&&u\ge\uRand{2}.
 \label{eq:rcu-B-at-star}
\end{align}

There is no missing branch from $A$.  On $2\le u\le3$, clearing the positive
denominators in $B_u(u)-A(u-1)$ and putting $v=u-2$ leaves
\[
 v^4+3v^3+3v^2+3v+3>0.
\]
On $3\le u\le\uRand{2}$, one has $B_u(u)>4/3$.  On
$\uRand{2}\le u\le\uRand{3}$, \eqref{eq:rcu-B-at-star} is at least $4/3$,
with equality exactly at $u=\uRand{3}$; afterwards $A(2)=4/3$ is larger.  Combining
these observations with \eqref{eq:rcu-binary-reduction} proves the announced
upper bound
\begin{equation}
 \frac{\min\{u,\mathcal T_u(D)\}}
      {1+\int_0^sS_D(t)^2\,dt}
 \le\RrandRO(u)
 \label{eq:rcu-announced-upper}
\end{equation}
for every finite distribution $D$.

The junction statements now follow directly.  Equations
\cref{eq:rcu-transition-cubic,eq:rcu-B-at-u,eq:rcu-B-at-star} give
continuity at $\uRand{2}$, and the square-root branch equals $4/3$ at
$u=\uRand{3}$.  Differentiating \eqref{eq:rcu-B-at-u} shows that its only
interior maximum above one occurs at $u=(3+\sqrt3)/2$; substitution gives
\eqref{eq:intro-random-common-maximum}.

\subsection{Worst-case upper bound from instance optimality}

No second learning argument is needed.  Apply the universal unannounced
policy of \cref{thm:common-upper-instance} to $D=D[M]$.  The endpoint
comparison \cref{eq:rcu-Phi-endpoint-upper} and the scalar bound
\cref{eq:rcu-announced-upper} give
\[
 \PhiRO(D[M])
 \le \RrandRO(u)\Omega_{\mathsf{RO},u}(D[M]).
\]
The exact empirical identity \cref{eq:rcu-offline-finite} has a nonnegative
linear correction.  Therefore the instance-optimal upper bound immediately
yields
\begin{equation}
 \EE\ALG_{\mathcal A_{n,u}}
 \le\RrandRO(u)\OPT_u+o_u(n^2),
 \label{eq:rcu-final-upper}
\end{equation}
which is \eqref{eq:intro-random-common-upper}.  For $0<u\le1$, raw execution
is exactly clairvoyantly optimal, so no learning is required.

\subsection{Binary oblivious lower bounds}

It remains to show that the two endpoint policies used above cannot be
improved in the worst case by a more adaptive algorithm.  We first record a
binary random-labeling envelope.

\begin{lemma}[Binary cap envelope]
\label{lem:rcu-binary-lower}
Fix $u>1$, $x\in[0,1]$, and put $z=1-x$.  Let $M_n$ have $xn+O(1)$ jobs of
value $u$ and $zn+O(1)$ zero jobs.  For every randomized nonanticipating
algorithm, even an announced one, some fixed labeling satisfies
\begin{equation}
 \frac{\EE\ALG}{n^2}\ge
 \begin{cases}
  \displaystyle\frac u2-o_u(1),&x\ge (u-1)/u,\\[6pt]
  \displaystyle\frac{1+x+ux^2}{2}-o_u(1),&x\le(u-1)/u.
 \end{cases}
 \label{eq:rcu-binary-fluid-lower}
\end{equation}
Every labeling has
\begin{equation}
 \frac{\OPT_u}{n^2}
 =\frac{1+(u-1)x^2}{2}+O_u(1/n).
 \label{eq:rcu-binary-offline}
\end{equation}
\end{lemma}

\begin{proof}
First fix a deterministic policy.  If it does not complete all jobs, the
claim is immediate.  Otherwise place the occurrences uniformly on the
public labels, distinguishing equal-valued occurrences as tokens.  By
deferred decisions, the analytically exposed occurrence word is a uniform
permutation, and the test/raw choice before its next entry is predictable.

Take $\delta_n=n^{-1/4}$.  The predictable-permutation bound
\cref{lem:predictable-permutation-sampling}, applied to the zero indicator
and the test selector, gives one event of probability $1-o(1)$ on which
every first-touch prefix
ending with at least $\delta_n n$ untouched labels and containing $t$ tests
satisfies
\begin{equation}
 \bigl|N_0(t)-zt\bigr|\le
 \Delta_n,
 \qquad
 \Delta_n=O(\delta_n^{-1}\sqrt{n\ln n})=o(n).
 \label{eq:rcu-predictable-urn}
\end{equation}
Here $N_0(t)$ is the number of tested zeros.  The event is simultaneous, so
we do not condition on the policy's adaptive stopping time.

Edit the suffix by running raw every still-untouched job that the original
schedule would later test, deleting its later known-processing operation.
Only $\delta_n n+O(1)$ jobs change, so the edit changes total completion cost
by $o_u(n^2)$.  Let $q$ be the final tested fraction of the edited schedule.
On \eqref{eq:rcu-predictable-urn}, every prefix of its completion curve is
dominated, up to $o(1)$ vertical error, by the divisible items
\begin{center}
\begin{tabular}{@{}lcc@{}}
\toprule
item & completion capacity & work per completion\\
\midrule
zero-test module & $zq$ & $1/z$\\
tested cap processing & $xq$ & $u$\\
raw jobs & $1-q$ & $u$\\
\bottomrule
\end{tabular}
\end{center}
When $z=0$, the zero-test module is simply omitted.
Indeed, $t$ tests complete at most $zt+\Delta_n$ zeros and expose at most
$xt+\Delta_n$ cap jobs; raw and known cap processing both use $u$ work per
completion.  The ordinary divisible-knapsack exchange therefore
upper-bounds completed mass at every work value.  Integrating the remaining
mass gives a lower bound on total completion time.

If $z\le1/u$, every displayed item costs at least $u$ per completion, so the
normalized area is at least $u/2-o_u(1)$.  If $z\ge1/u$, the zero-test module
comes first.  For fixed $q$ its remaining-mass area is
\begin{equation}
 F_z(q)=q-\frac{zq^2}{2}+\frac u2(1-zq)^2,
 \qquad
 F_z'(q)=(1-zq)(1-uz).
 \label{eq:rcu-binary-F}
\end{equation}
It is minimized at $q=1$, giving
$F_z(1)=(1+x+ux^2)/2$.  The suffix edit, vertical repair, and bad event all
contribute $o_u(n^2)$.

The argument holds for each deterministic policy.  For a randomized policy,
fix its seed, average the result over seeds and the finite uniform placement
space, and select a placement whose seed-expected cost is at least the joint
average.  This placement is fixed before the seed and is therefore an
oblivious input.  Finally, the effective lengths of the two job types are
$1$ and $u$, so \eqref{eq:rcu-binary-offline} follows from
\cref{lem:offline}.
\end{proof}

We now select the binary masses.  If $\uRand{1}<u\le\uRand{2}$, take
\begin{equation}
 x=\frac{u-1}{u}.
 \label{eq:rcu-lower-x-crossing}
\end{equation}
This is the boundary $z=1/u$ in \cref{lem:rcu-binary-lower}, and its online
and offline coefficients have ratio
\begin{equation}
 \frac{u}{1+(u-1)((u-1)/u)^2}
 =\frac{u^3}{u^2+(u-1)^3}.
 \label{eq:rcu-lower-first-branch}
\end{equation}
If $\uRand{2}\le u\le\uRand{3}$, take
\begin{equation}
 x=\frac1{\sqrt u-1}.
 \label{eq:rcu-lower-x-stationary}
\end{equation}
The definition of $\uRand{2}$ is precisely the point from which
$x\le(u-1)/u$.  Since $(u-1)x^2=1+2x$, the second line of
\eqref{eq:rcu-binary-fluid-lower} divided by
\eqref{eq:rcu-binary-offline} becomes
\begin{equation}
 1+\frac x2=1+\frac1{2(\sqrt u-1)}.
 \label{eq:rcu-lower-second-branch}
\end{equation}
Irrational masses are approximated by empirical rational frequencies; all
rounding and diagonal terms are $O_u(n)$.

For the last branch, the cap distribution is replaced by a fixed bounded
one.

\begin{lemma}[Balanced zero--two lower bound]
\label{lem:rcu-zero-two-lower}
For every $u\ge\uRand{3}$ and every randomized algorithm, arbitrarily large even
$n$ admit a fixed labeling with $n/2$ jobs of value zero and $n/2$ jobs of
value two such that
\begin{equation}
 \EE\ALG\ge n^2-o_u(n^2),
 \qquad
 \OPT_u=\frac34n^2+O(n).
 \label{eq:rcu-zero-two-costs}
\end{equation}
\end{lemma}

\begin{proof}
Use the same trace bijection, predictable-urn event, and suffix edit as in
\cref{lem:rcu-binary-lower}.  For a realized tested fraction $q$, the three
fluid items now have capacities and costs
\begin{center}
\begin{tabular}{@{}lcc@{}}
\toprule
item & completion capacity & work per completion\\
\midrule
zero-test module & $q/2$ & $2$\\
tested two processing & $q/2$ & $2$\\
raw jobs & $1-q$ & $u$\\
\bottomrule
\end{tabular}
\end{center}
Every item costs at least two per completion.  Thus the completion envelope
lies below $\min\{1,t/2\}$ independently of $q$, whose remaining-mass area is
one.  This proves the online estimate after the $o_u(n^2)$ repairs and finite
Yao selection \cref{lem:finite-yao}.  Since $u>3$, the two effective lengths
are one and three; their equally weighted minimum kernel gives offline
coefficient $3/4$.
\end{proof}

Equations \eqref{eq:rcu-lower-first-branch} and
\eqref{eq:rcu-lower-second-branch}, together with
\cref{lem:rcu-zero-two-lower}, match \eqref{eq:rcu-final-upper} in every
range.  For $0<u\le1$, raw execution is both an online and a clairvoyant
optimum.  This completes the proof of the exact curve in
\cref{thm:random-common-upper}.
